% ============================================================
%  IEEE Transactions — Kill-Chain MaxSAT Honeypot Placement
%  Student Writing Level Rewrite — v3
% ============================================================
\documentclass[journal]{IEEEtran}

\usepackage{amsmath,amssymb,amsfonts}
\usepackage{algorithmic}
\usepackage{algorithm}
\usepackage{array}
\usepackage{graphicx}
\usepackage{textcomp}
\usepackage{xcolor}
\usepackage{booktabs}
\usepackage{multirow}
\usepackage{url}
\usepackage{hyperref}
\usepackage{cite}
\usepackage{balance}
\usepackage{pifont}

\hypersetup{colorlinks=true,linkcolor=blue,citecolor=blue,urlcolor=blue}

\newcommand{\cmark}{\ding{51}}
\newcommand{\xmark}{\ding{55}}
\newcommand{\ie}{\emph{i.e.}}
\newcommand{\eg}{\emph{e.g.}}
\newcommand{\etal}{\emph{et~al.}}
\newcommand{\maxsat}{\textsc{MaxSAT}}
\newcommand{\rcII}{\textsc{RC2}}
\newcommand{\tildeW}{\tilde{w}}
\newcommand{\calK}{\mathcal{K}}
\newcommand{\calT}{\mathcal{T}}
\newcommand{\calA}{\mathcal{A}}
\newcommand{\calZ}{\mathcal{Z}}
\newcommand{\calG}{\mathcal{G}}
\newcommand{\calC}{\mathcal{C}}
\newcommand{\calPi}{\mathcal{\Pi}}

% ============================================================
\begin{document}

\title{Kill-Chain-Driven Honeypot Placement via\\
       Topology-Aware Weighted Partial \textsc{MaxSAT}:\\
       A Multi-Zone, Attack-Path Optimisation Framework\\
       for Enterprise Networks}

\author{
  \IEEEauthorblockN{Tu Nguyen (PhD Student) and Professor Haadi Jafarian}\\
  \IEEEauthorblockA{
    Department of Computer Science,
    University of Colorado Denver\\
    Denver, CO, USA\\
    \texttt{\{tu.nguyen, haadi.jafarian\}@ucdenver.edu}
  }
}

\markboth{IEEE TRANSACTIONS ON INFORMATION FORENSICS AND SECURITY}%
{Nguyen \& Jafarian: Kill-Chain-Driven Honeypot Placement via MaxSAT}

\maketitle

% ============================================================
\begin{abstract}
Deciding where to place honeypots in a large company network is
surprisingly hard. The network is split into different security
zones, contains thousands of different types of servers, and faces
attackers who follow well-known step-by-step attack patterns called
kill chains. On top of that, the organisation has a limited budget,
and some honeypot types cannot be deployed at the same time because
they would conflict with each other.

In this paper, we model honeypot placement as a \emph{Weighted
Partial Maximum Satisfiability} (\maxsat{}) problem. We encode six
goals at once: how well we detect threats, how many attack techniques
we cover from the MITRE ATT\&CK framework, how many attack tactic
families we cover, how well we stop attackers early, how well we
catch attackers at the end of their path, and how good our coverage
is across all attack paths. All hard rules---budget limits, zone
isolation, and conflict constraints---are encoded as clauses that
must never be violated.

We test our approach on networks ranging from 100 nodes (tiny) to
over 5.5 million nodes (very large), simulating five different
real-world attack scenarios. Our \maxsat{} solver achieves a
quality score of $Q = 67.4$ compared to $Q = 52.8$ for the best
competing approach ($+27.8\%$), while also providing a
mathematical guarantee that its answer is the best possible
solution.
\end{abstract}

\begin{IEEEkeywords}
Honeypot placement, Weighted Partial MaxSAT, kill chain,
MITRE ATT\&CK, network deception, multi-zone topology,
proactive cyber defence, RC2 solver, enterprise network security.
\end{IEEEkeywords}

% ============================================================
\section{Introduction}
\label{sec:intro}
% ============================================================

\IEEEPARstart{A}{ honeypot} is a fake computer system that is
designed to look attractive to attackers. When an attacker interacts
with it, the defender learns about their methods and can raise an
alert. Honeypots are widely used in network security, but their
usefulness depends heavily on \emph{where} they are placed in the
network. If a honeypot is put in the wrong place, it will never be
visited by attackers and the investment is wasted.

Real company networks make this placement decision much harder than
it sounds. Consider a typical large organisation. It has an
internet-facing area (called a DMZ) where its public web servers
sit. It has an internal office network. It might have cloud
services. It may also have industrial control systems (called
OT/SCADA) that run physical equipment like power plants or
factories. Each of these areas has different security rules,
different types of computers, and different risks. Attackers do not
wander randomly through these areas --- they follow structured
plans called \emph{kill chains}~\cite{hutchins2011intelligence},
moving step by step from one zone to another toward their final
target.

The problem is further complicated by three facts. First, the
organisation has a limited budget and cannot afford to deploy a
honeypot on every machine. Second, certain combinations of honeypot
types are incompatible --- for example, an industrial control system
honeypot deployed in the cloud zone would be useless because the
OT zone is physically isolated from the cloud. Third, there are
many possible attack paths through the network and an ideal
honeypot placement should intercept as many of them as possible.

Existing methods for honeypot placement do not handle all these
challenges at the same time. Simple rule-based approaches just
follow fixed guidelines. Greedy algorithms pick the best-looking
option at each step but cannot reverse a bad decision. Methods like
genetic algorithms can explore many solutions but have no way to
prove their answer is the best. In this paper, we take a different
approach: we use a formal solver that can consider \emph{all}
constraints and objectives together and provide a \emph{proven
optimal} answer.

Our main contributions are:

\begin{enumerate}

\item \textbf{A six-goal placement objective} encoded as a single
\maxsat{} optimisation problem. The six goals are: (1)
detection efficiency, (2) coverage of MITRE ATT\&CK attack
techniques, (3) coverage of MITRE ATT\&CK tactic families,
(4) catching attackers early in their path, (5) catching attackers
at the end of their path for forensic evidence, and (6) covering
as many attack paths as possible. Tactic-family coverage is
directly built into the solver as a bonus soft clause so the
solver actively tries to maximise it, not just measure it
afterwards.

\item \textbf{A priority system with four levels} for the
solver goals. The most important goals (catching attackers early)
are given weights 1000 times larger than the least important
goals (basic detection). This tells the solver to focus on
the most dangerous situations first. We also encode the key
insight that one honeypot can cover multiple attack paths at
once --- something no greedy method can discover.

\item \textbf{A thorough experimental evaluation} across six
different network sizes, using a Monte Carlo simulation to
model attacker behaviour and testing what happens when
threat intelligence updates cause us to re-run the solver
with new priorities.

\end{enumerate}

The rest of this paper is organised as follows.
Section~\ref{sec:related} discusses related work.
Section~\ref{sec:problem} defines the problem formally.
Section~\ref{sec:formulation} presents our \maxsat{} model.
Section~\ref{sec:methodology} describes our experiments.
Section~\ref{sec:results} presents the results.
Section~\ref{sec:dynamics} explains why our approach beats
the baselines.
Section~\ref{sec:discussion} covers limitations and future work.
Section~\ref{sec:conclusion} concludes.

% ============================================================
\section{Related Work}
\label{sec:related}
% ============================================================

\subsection{Honeypot Placement Research}

Honeypots have been studied for decades. Early work focused on
making individual honeypots realistic enough to fool
attackers~\cite{spitzner2003honeypots}. Provos built a system
that could run many virtual honeypots on one
machine~\cite{provos2004virtual}, which made large-scale deployment
practical. Spitzner popularised the idea of honeynets --- groups
of honeypots that together simulate an entire
network~\cite{spitzner2003honeypots}. However, most of this early
work assumed a human expert would decide where to put the honeypots.

More recently, researchers have tried to automate the placement
decision. Some treated it as a game between an attacker and a
defender~\cite{yuill2004honeypots}. Tan \etal{} and Pawlick
\etal{} used game theory to model how attackers respond to
deception~\cite{tan2011,pawlick2019game}. These are interesting
approaches but they typically assume a simple two-player game on a
flat network. In reality, enterprise networks have multiple zones
with different trust levels, and attackers can follow many different
simultaneous paths.

\subsection{Using MITRE ATT\&CK for Defence Planning}

MITRE ATT\&CK~\cite{mitre2023} is a publicly available catalogue
of the tactics and techniques that real attackers use. It has
become the standard reference for both attackers (red teams) and
defenders (blue teams). Applebaum \etal{} showed how it could be
used to automatically generate realistic attack
simulations~\cite{applebaum2016intelligent}. Hemberg \etal{} used
evolutionary algorithms to select the best defensive controls from
the ATT\&CK catalogue~\cite{hemberg2020}. However, their approach
did not consider budget limits or network zone topology. Our work
directly embeds ATT\&CK technique weights, stealth scores, and
tactic-family coverage into the solver as weighted constraints.

\subsection{Using Formal Methods for Network Security}

Formal methods use mathematical logic to prove properties of
systems. In network security, they have been used to check
firewall rules~\cite{al2011firewall}, automatically generate
network configurations~\cite{jeffrey2009}, and verify security
policies~\cite{fisler2005}. Closer to our work, Pinisetty \etal{}
applied \maxsat{} to choose the best rules for an intrusion
detection system~\cite{pinisetty2017}. However, to our knowledge,
no prior work has applied Weighted Partial \maxsat{} to the
honeypot placement problem with kill-chain-aware objectives and
physical zone isolation constraints.

Integer Linear Programming (ILP) has also been used for security
resource allocation~\cite{saha2015optimal,bozorgi2010}. ILP can
find optimal solutions but becomes very slow as the problem grows.
Our zone-decomposed approach keeps the solver fast by splitting
the large network into smaller per-zone sub-problems.

\subsection{Moving-Target Defence}

Moving-target defence (MTD)~\cite{zhu2013game} is a strategy where
the defender constantly changes the network configuration to confuse
attackers. This is a complementary idea to honeypots: MTD changes
the real system to create uncertainty, while honeypots add fake
systems to create traps. Our formulation focuses on placing the
best traps, not on moving real assets.

\subsection{How Our Work Compares}

Table~\ref{tab:related} compares our approach with the most
closely related prior work across five key dimensions. Our method
is the only one that satisfies all five requirements.

\begin{table}[htbp]
\caption{Comparison of Honeypot Placement Approaches}
\label{tab:related}
\centering
\renewcommand{\arraystretch}{1.2}
\begin{tabular}{@{}lccccc@{}}
\toprule
\textbf{Approach} &
\rotatebox{60}{\textbf{Multi-Zone}} &
\rotatebox{60}{\textbf{Kill-Chain}} &
\rotatebox{60}{\textbf{ATT\&CK}} &
\rotatebox{60}{\textbf{Budget+Conflict}} &
\rotatebox{60}{\textbf{Optimal}} \\
\midrule
Game-theoretic \cite{tan2011}    & \xmark & \xmark & \xmark & Partial & Nash eq. \\
GA-based \cite{hemberg2020}      & \xmark & \xmark & \cmark & \xmark  & None \\
ILP sensor \cite{saha2015optimal}& \cmark & \xmark & \xmark & \cmark  & Yes (ILP) \\
Greedy-BiDir (our baseline)      & \cmark & \cmark & Partial& \xmark  & None \\
\textbf{Ours (MaxSAT)}           & \cmark & \cmark & \cmark & \cmark  & \textbf{Proven} \\
\bottomrule
\end{tabular}
\end{table}

% ============================================================
\section{Problem Definition}
\label{sec:problem}
% ============================================================

\subsection{What We Are Trying to Solve}

At a high level, we are trying to answer: \emph{``Given a limited
budget and a set of honeypot types, which honeypots should we deploy
and where, so that we catch the most attackers across all possible
attack paths?''}

We formally define the problem instance as the tuple:
\begin{equation}
  \mathcal{I} =
  (\calK,\calT,\calA,\calZ,\calG,\calC,\mathcal{I}_z,
   \calPi,w,\mathrm{cost},B,B_z,dm,hd,\sigma,\rho,iv)
\end{equation}

Each symbol in this tuple represents a different part of the
problem, explained in the following sections.

\subsection{The Basic Building Blocks (Sets)}

\begin{table}[htbp]
\caption{Problem Sets}
\label{tab:sets}
\centering
\renewcommand{\arraystretch}{1.2}
\begin{tabular}{@{}ll@{}}
\toprule
\textbf{Symbol} & \textbf{What It Means} \\
\midrule
$\calK$ & The list of honeypot types we can deploy \\
$\calT$ & All the attack techniques from the MITRE ATT\&CK list \\
$\calA$ & All the individual computer assets in the network \\
$\calZ$ & The security zones (DMZ, Internal LAN, Cloud, etc.) \\
$\calG$ & The types of servers (web server, SSH, database, etc.) \\
$\calPi$ & The attack paths (sequences of steps an attacker takes) \\
$\calC \subseteq \calK \times \calK$ & Pairs of honeypots that cannot be deployed together \\
$\mathcal{I}_z \subseteq \calZ \times \calZ$ & Pairs of zones that are physically isolated (air gaps) \\
\bottomrule
\end{tabular}
\end{table}

\subsection{The Numbers We Need (Parameters)}

\begin{table}[htbp]
\caption{Problem Parameters}
\label{tab:params}
\centering
\renewcommand{\arraystretch}{1.2}
\begin{tabular}{@{}lll@{}}
\toprule
\textbf{Symbol} & \textbf{Type} & \textbf{What It Means} \\
\midrule
$ZK_i \subseteq \calZ$        & per honeypot & Which zones this honeypot can go in \\
$GK_i \subseteq \calG$        & per honeypot & Which server types this honeypot mimics \\
$S_{i,a} \subseteq \calT$     & per honeypot/asset & Which attacks it can detect on asset $a$ \\
$w_{j,a} \in \mathbb{R}^+$    & per attack/asset & How valuable detecting this attack is \\
$\mathrm{cost}_i \in \mathbb{R}^+$ & per honeypot & How much it costs to deploy \\
$B \in \mathbb{R}^+$          & global & Our total budget \\
$B_z \in \mathbb{R}^+$        & per zone & Budget limit for each zone \\
$dm_a \in \mathbb{R}^+$       & per asset & How important this asset's role is \\
$hd_a \in \mathbb{Z}^+$       & per asset & How many hops from the attacker's entry point \\
$\sigma_j \in [0,1]$          & per attack & How sneaky this attack is (0 = obvious, 1 = very stealthy) \\
$\rho_\pi \in [0,1]$          & per path & How likely attackers are to use this path \\
$iv_{\pi,h} \in \mathbb{R}^+$ & per path/hop & How valuable it is to catch the attacker at this step \\
\bottomrule
\end{tabular}
\end{table}

\subsection{Calculated Values}

Before running the solver, we compute four derived values that
combine the raw parameters into more meaningful weights.

\subsubsection{Which Assets Can a Honeypot Reach?}
Each honeypot type $k_i$ can only be deployed on certain server
types in certain zones. The set of assets it can cover is:
\begin{equation}
  T^*(i) = DK_i \;\cup\;
  \bigl\{\,a \in \calA :
    \mathrm{type}(a) \in GK_i \;\wedge\; \mathrm{zone}(a) \in ZK_i
  \bigr\}
\end{equation}

\subsubsection{Can This Honeypot Detect This Attack on This Asset?}
A simple yes/no indicator:
\begin{equation}
  R[j,i,a] =
  \begin{cases}
    1 & \text{if the honeypot can reach asset $a$ AND detect attack $t_j$ there} \\
    0 & \text{otherwise}
  \end{cases}
\end{equation}

\subsubsection{Adjusted Detection Value (Topology Weight)}
Not all detections are equally important. Catching an attacker
near the entry point is more valuable than catching them deep
inside the network (because we can stop them sooner). We scale
the raw value by the asset's role importance and divide by how
many hops away it is:
\begin{equation}
  \tildeW_{j,a} = w_{j,a} \times dm_a \times \frac{1}{hd_a}
\end{equation}

\textit{Example:} A firewall gateway asset with $dm=1.3$ and
$hd=1$ (right at the entry) gives $\tildeW = 1.30 \times w$.
A deep database with $dm=1.5$ but $hd=3$ (three hops deep)
gives $\tildeW = 0.50 \times w$ --- actually worth less to cover,
even though the database is more valuable, because catching the
attacker at the gate stops more damage.

\subsubsection{Stealth-Adjusted Weight}
Stealthy attacks are harder to detect in practice. If we do catch
them, it is a bigger win. We multiply by $(1 + \sigma_j)$ to give
stealthy techniques higher weight:
\begin{equation}
  W_{j,a} = \tildeW_{j,a} \times (1 + \sigma_j)
\end{equation}

\textit{Example:} T1048 Exfiltration is very stealthy
($\sigma=0.9$). With $w=2.0$ on a gateway:
$W = 2.0\times1.30\times1.90 = 4.94$.
T1046 Network Scanning is very obvious ($\sigma=0.1$).
With $w=0.5$: $W = 0.5\times1.00\times1.10 = 0.55$.
This means detecting exfiltration is worth about 9$\times$ more
than detecting a scan.

\subsubsection{Path Intercept Weight}
Finally, we combine the detection value with the attack path's
probability and how urgent it is to catch the attacker at that
step:
\begin{equation}
  PW_{\pi,h,j,a} = \rho_\pi \times iv_{\pi,h} \times W_{j,a}
\end{equation}

This means high-probability paths get heavier solver weights,
so the solver naturally prioritises the most likely attack
scenarios first.

\subsection{What the Solver Decides (Decision Variables)}

The solver needs to make four types of decisions, each
represented by a binary yes/no variable:

\begin{table}[htbp]
\caption{Decision Variables}
\label{tab:vars}
\centering
\renewcommand{\arraystretch}{1.2}
\begin{tabular}{@{}lll@{}}
\toprule
\textbf{Variable} & \textbf{Yes or No?} & \textbf{What It Decides} \\
\midrule
$x_i$      & $\{0,1\}$ & Should we deploy honeypot type $k_i$? \\
$c_{j,a}$  & $\{0,1\}$ & Are we detecting attack $t_j$ on asset $a$? \\
$p_{\pi,h}$& $\{0,1\}$ & Are we blocking path $\pi$ at step $h$? \\
$e_\pi$    & $\{0,1\}$ & Are we catching the attacker \emph{early} on path $\pi$? \\
\bottomrule
\end{tabular}
\end{table}

% ============================================================
\section{Our MaxSAT Model}
\label{sec:formulation}
% ============================================================

\maxsat{} works by encoding a problem as a set of logical clauses.
\emph{Hard clauses} are rules that must \emph{always} be obeyed.
\emph{Soft clauses} are goals we \emph{want} to achieve as much as
possible, each with a weight representing how important they are.
The solver finds the assignment of variables that satisfies all
hard clauses and maximises the total weight of satisfied soft
clauses.

\subsection{Hard Clauses: Rules That Cannot Be Broken}

We have seven types of hard constraints.

\paragraph{C1 --- You Can Only Claim a Detection If a Honeypot Is Actually Deployed There}
We cannot count a detection unless at least one deployed honeypot
covers that (attack, asset) pair:
\begin{equation}
  \neg c_{j,a} \;\vee\;
  \bigl(x_{i_1} \vee x_{i_2} \vee \cdots \vee x_{i_k}\bigr)
\end{equation}
In plain English: ``If we say we are detecting attack $t_j$ on asset
$a$, then at least one of the honeypots that can catch $t_j$ on $a$
must be deployed.''

\paragraph{C2 --- Do Not Exceed the Total Budget}
\begin{equation}
  \sum_{i \in \calK} \mathrm{cost}_i \cdot x_i \;\leq\; B
\end{equation}

\paragraph{C3 --- Do Not Exceed Each Zone's Budget}
Each zone has its own budget slice. We cannot overspend in any
single zone:
\begin{equation}
  \sum_{\substack{i \in \calK \\ ZK_i \cap \{z\} \neq \emptyset}}
  \!\!\mathrm{cost}_i \cdot x_i \;\leq\; B_z
  \quad\forall\, z \in \calZ
\end{equation}

\paragraph{C4 --- Do Not Deploy Conflicting Honeypots}
Some honeypot types cannot coexist. If they conflict, only one can
be deployed:
\begin{equation}
  \neg x_i \;\vee\; \neg x_l
  \quad \forall\,(k_i,k_l) \in \calC
\end{equation}

Our model has seven conflict pairs:
\texttt{scada\_trap} vs \texttt{web\_trap},
\texttt{scada\_trap} vs \texttt{db\_trap},
\texttt{scada\_trap} vs \texttt{ssh\_trap},
\texttt{scada\_trap} vs \texttt{dns\_trap},
\texttt{generic\_trap} vs \texttt{ad\_trap},
\texttt{ad\_trap} vs \texttt{web\_trap},
\texttt{smb\_trap} vs \texttt{dns\_trap}.
These conflict rules also help the solver run faster by eliminating
impossible combinations early.

\paragraph{C5 --- Respect Physical Zone Isolation (Air Gaps)}
Some zones are physically disconnected. No honeypot can pretend to
span both sides of an air gap:
\begin{equation}
  \forall\,(z_a,z_b) \in \mathcal{I}_z,\;
  \forall\, i \in \calK :
  z_a \in ZK_i \wedge z_b \in ZK_i
  \;\Longrightarrow\; x_i = 0
\end{equation}

\paragraph{C6 --- A Path Hop Is Only Covered If We Actually Detect Something There}
We cannot claim we are blocking an attack path at step $h$ unless
at least one technique used at that step is detected on at least
one asset in that zone:
\begin{equation}
  \neg p_{\pi,h} \;\vee\;
  \bigl(c_{j_1,a_1} \vee \cdots \vee c_{j_n,a_n}\bigr)
\end{equation}

\paragraph{C7 --- Early Interception Must Actually Happen Before the Final Step}
We can only claim early interception of path $\pi$ if the
attacker is caught before reaching their final target:
\begin{equation}
  e_\pi = 1 \;\Longrightarrow\;
  \exists\, h < |\pi| : p_{\pi,h} = 1
\end{equation}

\subsection{Soft Clauses: Goals We Want to Maximise}

Soft clauses represent what we \emph{want} the solver to achieve.
We use four priority levels. Think of these like four buckets of
goals, where the top bucket is 1000 times more important than the
second bucket, which is 100 times more important than the third,
and so on. The solver fills the most important bucket first.

\subsubsection{Level 4 (Most Important): Catch Attackers Early}
If we can catch an attacker before they reach their final target,
we prevent the damage entirely. This is much better than just
detecting them after the damage is done. We give this goal the
highest weight:
\begin{align}
  \text{weight} &= \rho_\pi \times
    \max_{h < |\pi|}(iv_{\pi,h}) \times
    \sum_{j,a} W_{j,a} \notag \\
  \text{goal:} &\quad (e_\pi) \quad \text{for every attack path } \pi
\end{align}

\subsubsection{Level 3: Cover As Many Attack Path Steps As Possible}

\textbf{Forward coverage} (catching the attacker as they move
toward their target):
\begin{align}
  \text{weight} &= \rho_\pi \cdot iv_{\pi,h} \cdot W_{j,a} \notag \\
  \text{goal:} &\quad (p_{\pi,h})
    \quad \text{for every path } \pi \text{ and step } h
\end{align}

\textbf{Backward coverage} (catching the attacker during or after
their final action --- useful for forensics and evidence, but less
important than preventing the attack in the first place):
\begin{align}
  \text{weight} &= \rho_\pi \cdot iv_{\pi,h} \cdot W_{j,a} \times 0.7 \notag \\
  \text{goal:} &\quad (p_{\pi,h})
    \quad \text{for steps in reverse order}
\end{align}

The 0.7 factor means forward detection is always prioritised over
backward detection, reflecting the real-world importance of
prevention over forensics.

\noindent\textbf{Key insight about cross-path overlap:}
Because we encode \emph{all} attack paths in the same solver
instance, a single honeypot can satisfy goals for multiple paths
at once. For example, an SSH honeypot in the internal LAN zone
that detects technique T1021 simultaneously covers the
\textsc{web-to-db}, \textsc{ot-infiltration}, and
\textsc{brute-to-ad} paths --- all three at once. A greedy
algorithm would only count this honeypot's contribution toward
whichever path it is currently evaluating and miss the other two.
Our solver discovers this ``force multiplier'' effect automatically.

\subsubsection{Level 2: Cover as Many Attack Techniques and Tactic Families as Possible}

For each individual MITRE ATT\&CK technique, we want to detect it
at least once somewhere in the network:
\begin{align}
  \text{weight} &= \sum_{a \in \calA} W_{j,a} \notag \\
  \text{goal:} &\quad
    \Bigl(\text{at least one honeypot detects technique } t_j\Bigr)
    \quad \forall\, j \in \calT
  \label{eq:tier2-tech}
\end{align}

We also want to cover all MITRE tactic \emph{families} (groups of
related techniques like ``Initial Access'' or ``Lateral Movement'').
A $1.2\times$ bonus is given for covering at least one technique
per family:
\begin{align}
  \text{weight} &= 1.2 \times \sum_{j \in f} \sum_{a} W_{j,a}
  \notag \\
  \text{goal:} &\quad
    \Bigl(\text{at least one technique in family } f \text{ is detected}\Bigr)
    \quad \forall\, f \in \mathcal{F}
  \label{eq:tier2-fam}
\end{align}

The bonus means the solver prefers to cover \emph{one technique
from each of ten families} over \emph{ten techniques from the same
family}, which gives much better overall ATT\&CK coverage.

\subsubsection{Level 1 (Least Important): Basic Detection Opportunities}
At the lowest priority, we simply want each individual
(technique, asset) detection pair to be covered:
\begin{equation}
  \text{weight} = W_{j,a},
  \quad \text{goal: detect } t_j \text{ on } a
  \quad \forall\, j, a
\end{equation}

\subsection{The Complete Optimisation Objective}

Putting it all together:
\begin{equation}
\max \sum_{\text{Level 4}} w_4 \cdot \text{sat}
   + \sum_{\text{Level 3}} w_3 \cdot \text{sat}
   + \sum_{\text{Level 2}} w_2 \cdot \text{sat}
   + \sum_{\text{Level 1}} w_1 \cdot \text{sat}
\end{equation}
subject to all hard constraints C1--C7, where
$w_4 = 1000\,w_3$, $w_3 = 100\,w_2$, $w_2 = 10\,w_1$.

\subsection{How We Measure the Final Answer (Quality Score $Q$)}

After the solver finishes, we compute a single quality score $Q$
that combines five percentages:
\begin{equation}
  Q = 0.35\cdot\text{DetEff} + 0.25\cdot\text{TechCov}
    + 0.15\cdot\text{FamCov} + 0.15\cdot\text{FwdPath}
    + 0.10\cdot\text{BwdPath}
\end{equation}

Each component is a percentage from 0 to 100:

\begin{align}
  \text{DetEff}\%   &= \frac{\text{Total weighted detections achieved}}
                             {\text{Total possible weighted detections}} \times 100 \\[4pt]
  \text{TechCov}\%  &= \frac{\text{Number of ATT\&CK techniques detected}}
                             {\text{Total number of techniques}} \times 100 \\[4pt]
  \text{FamCov}\%   &= \frac{\text{Number of ATT\&CK tactic families covered}}
                             {\text{Total tactic families}} \times 100 \\[4pt]
  \text{FwdPath}\%  &= \frac{\text{Number of (path, step) pairs covered forward}}
                             {\text{Total (path, step) pairs}} \times 100 \\[4pt]
  \text{BwdPath}\%  &= \frac{\text{Number of (path, step) pairs covered backward}}
                             {\text{Total (path, step) pairs}} \times 100
\end{align}

\subsection{Splitting the Problem by Zone to Speed It Up}

For large networks, we decompose the problem by zone. Instead of
solving one giant problem, we solve one smaller problem per zone:
\begin{equation}
  \forall\, z \in \calZ:\;
  \text{solve } \mathcal{I}_z = (\calK_z, \calT_z, \calA_z, \calC_z, B_z)
\end{equation}
where $\calK_z$ contains only the honeypot types that can operate
in zone $z$. Each zone problem runs in parallel on a separate CPU
core, and a final budget check ensures we did not overspend
across zones.

\subsection{Updating the Placement When Threat Intelligence Changes}
\label{sec:redeployment}

In the real world, the threat landscape changes. A new vulnerability
report might make one attack path more likely. When this happens,
we need to update our honeypot deployment. Algorithm~\ref{alg:redeploy}
describes how we do this efficiently using the previous solution as
a starting point.

\begin{algorithm}[htbp]
\caption{Adaptive Redeployment When Threat Priors Change}
\label{alg:redeploy}
\begin{algorithmic}[1]
\REQUIRE Old deployment $x^*$; new attack path probabilities $\rho'_\pi$
\ENSURE  New deployment $x'^*$ that is optimal for the updated threat model
\STATE Recalculate path weights: $PW'_{\pi,h,j,a} \leftarrow \rho'_\pi \cdot iv_{\pi,h} \cdot W_{j,a}$
\STATE Update the soft clause weights in the solver (hard rules C1--C7 stay the same)
\STATE Start the solver from the old solution $x^*$ (warm start --- much faster than from scratch)
\STATE Re-run the solver to get new solution $x'^*$
\STATE Calculate improvement: $\Delta = Q(x'^*) - Q(x^*)$
\IF{$\Delta > 0$}
  \STATE Use the new deployment $x'^*$ and record the improvement
\ELSE
  \STATE Keep the old deployment $x^*$ (it is still optimal)
\ENDIF
\RETURN $x'^*$
\end{algorithmic}
\end{algorithm}

The warm start in Step 3 is important: since only the soft clause
weights changed (not the hard rules), the solver can reuse most of
its previous work and converge much faster than solving from scratch.

% ============================================================
\section{Experimental Setup}
\label{sec:methodology}
% ============================================================

\subsection{The Network We Tested On}

We built a synthetic enterprise network with five security zones.
Table~\ref{tab:zones} shows the zones and their properties. Two
zones are physically isolated via air gaps: the OT/SCADA zone
(industrial equipment) and the Management zone cannot connect to
the internet-facing DMZ or the cloud.

\begin{table}[htbp]
\caption{Network Zones}
\label{tab:zones}
\centering
\renewcommand{\arraystretch}{1.2}
\begin{tabular}{@{}llccc@{}}
\toprule
\textbf{Zone} & \textbf{Description} &
\textbf{Trust Level} & \textbf{\% of Assets} & \textbf{\% of Attacks} \\
\midrule
zone1 & Internet-Facing / DMZ    & Low (0)       & 15\% & 45\% \\
zone2 & Internal Company Network & Medium (2)    & 40\% & 25\% \\
zone3 & Cloud / Hybrid           & Low-Med (1)   & 25\% & 20\% \\
zone4 & OT / Industrial Systems  & High (3)      & 10\% &  5\% \\
zone5 & Management Network       & Highest (4)   & 10\% &  5\% \\
\bottomrule
\end{tabular}
\end{table}

The network has nine connections between zones, with some acting
as chokepoints --- single entry/exit points where all traffic must
pass. These chokepoints are especially important because a honeypot
near a chokepoint can intercept every attacker that uses that route.

\subsection{The Attack Scenarios We Modelled}

We modelled five realistic attack paths that an attacker might
follow (Table~\ref{tab:paths}). Each path has a probability
($\rho$) representing how likely real attackers are to use it.

\begin{table}[htbp]
\caption{The Five Attack Scenarios}
\label{tab:paths}
\centering
\renewcommand{\arraystretch}{1.2}
\begin{tabular}{@{}llccc@{}}
\toprule
\textbf{Name} & \textbf{What the Attacker Does} & \textbf{Prob.} & \textbf{Starts} & \textbf{Target} \\
\midrule
web\_to\_db      & Hack website, steal database   & 0.30 & zone1 & zone2 \\
cloud\_pivot     & Use cloud to enter internal net & 0.25 & zone3 & zone2 \\
brute\_to\_ad    & Guess passwords, take over AD  & 0.20 & zone1 & zone5 \\
ot\_infiltration & Sabotage industrial systems    & 0.15 & zone1 & zone4 \\
ransomware       & Email phishing, spread malware & 0.10 & zone2 & zone2 \\
\bottomrule
\end{tabular}
\end{table}

\subsection{Network Sizes We Tested}

We ran experiments at six different scales to see how well our
approach holds up as the network grows.

\begin{table}[htbp]
\caption{Network Sizes Tested}
\label{tab:scales}
\centering
\renewcommand{\arraystretch}{1.2}
\begin{tabular}{@{}lrrccc@{}}
\toprule
\textbf{Size} & \textbf{Computers} & \textbf{Budget} &
\textbf{Zones} & \textbf{Attacks} & \textbf{Honeypot Types} \\
\midrule
Tiny    &          100 &    250 & 2 & 2 & 3 \\
Small   &          500 &    500 & 3 & 3 & 4 \\
Medium  &        5,000 &  1,750 & 4 & 3 & 5 \\
Large   &       50,000 & 12,500 & 5 & 5 & 7 \\
XLarge  &      500,000 & 62,500 & 5 & 5 & 8 \\
XXLarge &    5,560,000 &128,000 & 5 & 5 & 8 \\
\bottomrule
\end{tabular}
\end{table}

\subsection{The Honeypot Types Available}

We defined eight types of honeypots, each designed to attract a
specific kind of attack (Table~\ref{tab:profiles}).

\begin{table}[htbp]
\caption{Available Honeypot Types}
\label{tab:profiles}
\centering
\renewcommand{\arraystretch}{1.2}
\begin{tabular}{@{}llcc@{}}
\toprule
\textbf{Honeypot} & \textbf{Goes In} & \textbf{Cost} & \textbf{Detects} \\
\midrule
web\_trap     & DMZ, Cloud          & 1.0 & T1190, T1133, T1059 \\
ssh\_trap     & DMZ, Internal, Cloud& 0.8 & T1110, T1021, T1078 \\
db\_trap      & Internal, Cloud     & 1.3 & T1213, T1048, T1485 \\
smb\_trap     & Internal only       & 0.9 & T1021, T1550, T1486 \\
scada\_trap   & OT zone only        & 2.0 & T0855, T0814, T0869 \\
ad\_trap      & Internal, Mgmt      & 1.5 & T1558, T1003, T1098 \\
dns\_trap     & DMZ, Internal, Cloud& 0.7 & T1572, T1095, T1046 \\
generic\_trap & Anywhere            & 0.5 & T1046, T1082, T1110 \\
\bottomrule
\end{tabular}
\end{table}

\subsection{The Comparison Methods (Baselines)}

We compared our \maxsat{} approach against three simpler methods:

\begin{itemize}
  \item \textbf{Random}: Pick honeypots randomly until we run out
        of budget. We do this 100 times and take the average.
  \item \textbf{Greedy-Value}: At each step, pick the honeypot that
        gives the most detection value per dollar. Stop when the
        budget runs out.
  \item \textbf{Greedy-BiDir}: Like Greedy-Value but also considers
        attack path direction --- gives slightly less credit for
        backward path coverage (0.7 discount). This is the
        strongest baseline.
\end{itemize}

\subsection{Solver Settings}

We used the \rcII{} solver from the pysat Python library with a
60-second time limit. Each zone gets its own solver instance with
12 to 30 seconds, running in parallel on up to 4 CPU cores.

\subsection{Monte Carlo Simulation: Testing With Fake Attackers}
\label{sec:montecarlo}

To double-check that our formal model matches reality, we ran
a separate simulation using random attackers (Table~\ref{tab:montecarlo}).

\begin{table}[htbp]
\caption{How Many Simulated Attacks We Ran}
\label{tab:montecarlo}
\centering
\renewcommand{\arraystretch}{1.2}
\begin{tabular}{@{}lrr@{}}
\toprule
\textbf{Network Size} & \textbf{Simulated Attacks} & \textbf{Random Seed} \\
\midrule
Tiny    &       5,000 & 42 \\
Small   &      10,000 & 42 \\
Medium  &      50,000 & 42 \\
Large   &     500,000 & 42 \\
XLarge  &   5,000,000 & 42 \\
XXLarge &  25,600,000 & 42 \\
\bottomrule
\end{tabular}
\end{table}

In each simulated attack, we randomly pick one of the five attack
paths according to its probability, then step through the path
hop by hop. At each step, the attacker is caught if a deployed
honeypot covers that zone and technique. We measure three things:
(1) what percentage of attackers were caught early (before the
final step), (2) the average number of steps before being caught,
and (3) what percentage of attackers got through without being
caught at all.

We also ran 1,000 trials where we slightly changed the attack
path probabilities and then re-ran Algorithm~\ref{alg:redeploy}
to see how much the placement improved.

% ============================================================
\section{Results}
\label{sec:results}
% ============================================================

\subsection{Large Network (500,000 Computers)}

Table~\ref{tab:xlarge} shows how our method compares to the best
baseline on the XLarge network. The solver picked 4 honeypot
configurations using 98\% of the budget.

\begin{table}[htbp]
\caption{XLarge Results: Our Method vs. Best Baseline}
\label{tab:xlarge}
\centering
\renewcommand{\arraystretch}{1.2}
\begin{tabular}{@{}lrrrcc@{}}
\toprule
\textbf{Metric} & \textbf{Ours} & \textbf{Best Baseline} &
\textbf{Difference} & \textbf{\% Better} & \textbf{Over 10\%?} \\
\midrule
Detection Efficiency & 59.4  & 57.3  & +2.1  & +3.6\%   & \cmark \\
Technique Coverage   & 63.2  & 31.6  & +31.6 & +100.0\% & \cmark \\
Family Coverage      & 66.7  & 41.7  & +25.0 & +60.0\%  & \cmark \\
Forward Path Cover   & 73.3  & 60.0  & +13.3 & +22.2\%  & \cmark \\
Backward Path Cover  & 73.3  & 60.0  & +13.3 & +22.2\%  & \cmark \\
Early Intercept \%   & 100.0 & 100.0 & +0.0  & +0.0\%   & ---    \\
\midrule
\textbf{Q Score}     & \textbf{67.4} & \textbf{52.8} &
  \textbf{+14.7} & \textbf{+27.8\%} & \cmark \\
\bottomrule
\end{tabular}
\end{table}

The most striking result is Technique Coverage: our method covers
\emph{twice} as many MITRE ATT\&CK techniques as the best baseline
($+100\%$). This is because our solver simultaneously considers all
technique coverage objectives, while greedy methods only optimise
one step at a time and miss the cross-path overlap effect.

\subsection{Medium Network (5,000 Computers)}

At Medium scale, our solver achieved $Q = 91.91$ --- near perfect
--- using only 10.9\% of the available budget. All MITRE tactic
families were covered (FamCov = 100\%). The Monte Carlo simulation
confirmed that 100\% of simulated attackers were caught early, with
a 0\% miss rate, validating our formal model.

\subsection{How Performance Changes With Network Size}

\begin{table}[htbp]
\caption{Quality Score at Each Network Size}
\label{tab:scale}
\centering
\renewcommand{\arraystretch}{1.2}
\begin{tabular}{@{}lrrrr@{}}
\toprule
\textbf{Size} & \textbf{Our Score} & \textbf{Best Baseline} &
\textbf{Improvement} & \textbf{Honeypots Picked} \\
\midrule
Tiny    & 95.2 & 93.1 &  +2.1 & 2 of 3 \\
Small   & 92.4 & 88.7 &  +3.7 & 3 of 4 \\
Medium  & 91.9 & 84.2 &  +7.7 & 4 of 5 \\
Large   & 78.6 & 61.3 & +17.3 & 5 of 7 \\
XLarge  & 67.4 & 52.8 & +14.6 & 4 of 8 \\
XXLarge & 58.3 & 44.1 & +14.2 & 4 of 8 \\
\bottomrule
\end{tabular}
\end{table}

An important trend: the advantage of our method \emph{grows} as
the network gets larger. At Tiny scale, the improvement is only
$+2.1$ points because the problem is simple enough that even
random selection does well. At Large and XLarge scale, where the
problem becomes genuinely hard, our method outperforms the best
baseline by over $+14$ points consistently.

% ============================================================
\section{Why Our Method Beats the Baselines}
\label{sec:dynamics}
% ============================================================

\subsection{The Fundamental Problem With Greedy Methods}

All three baselines treat forward detection (catching the attacker
as they advance) and backward detection (catching them at the end)
as the same thing. They are not.

\textbf{Forward detection = prevention.} If we catch an attacker
at the first step of the web-to-database attack, the database is
never compromised. The company's data is safe.

\textbf{Backward detection = forensics.} If we only catch the
attacker at the very last step, the data has already been stolen.
We have evidence, but no prevention.

Greedy methods apply the same value to both, which leads to
suboptimal placements.

\subsection{Five Ways Our Approach Is Fundamentally Better}

\textbf{1 --- We prioritise early catches.}
Each step of an attack path has an ``intercept value'' that
increases toward the final step (for web-to-db: $1.5 \to 1.8
\to 2.0$). Our Level-4 goals explicitly reward catching the
attacker at any non-final step. Greedy methods see a higher
individual value at the final step and tend to focus there.

\textbf{2 --- We exploit overlap between attack paths.}
Techniques T1021 and T1078 appear in the \emph{same zone} on
three different attack paths simultaneously. Deploying an SSH
honeypot there covers all three paths at once. A greedy method
credits this honeypot toward only the one path it is currently
evaluating. Our solver encodes all paths simultaneously and
automatically discovers this $3\times$ multiplier.

\textbf{3 --- We reason about consequences of choices.}
If deploying the SCADA honeypot would cause three other honeypots
to be removed (due to conflict rules), is it worth it? Our solver
can reason about this full consequence chain. Greedy methods
cannot look ahead --- they just pick the best-looking option at
each step and cannot undo bad decisions.

\textbf{4 --- We understand the difference between trap types.}
Forward-path honeypots should look convincingly real (high
fidelity) to intercept the attacker without triggering evasion.
Backward-path honeypots at exfiltration points serve a different
purpose: confirming that data was stolen. Our 0.7 discount
on backward coverage encodes this two-layer deception strategy
directly. No baseline makes this distinction.

\textbf{5 --- We can respond optimally to new threats.}
When threat intelligence changes (say, a new CVE makes the
ransomware path three times more likely), our solver can be
re-run with the updated probabilities and will produce a
\emph{provably optimal} updated deployment in a fraction of
the original solve time. Baselines would need to re-run from
scratch with no guarantee of improvement.

\subsection{How Much Would Pure Path Improvements Help?}

If we only applied path-related improvements (better zone budgets,
soft conflict handling, attack-share weighting), we project:
FwdPath\% would rise from 73.3\% to 88--93\%;
BwdPath\% would rise from 73.3\% to 78--85\%;
and $Q$ would improve by $+3.25$ points from path dynamics alone.
Combined with all other improvements, we project $Q$ reaching
88--95, a 23--33\% gain over the best baseline.

% ============================================================
\section{Limitations and Future Work}
\label{sec:discussion}
% ============================================================

\subsection{What Our Approach Cannot Do Yet}

\textbf{We used a synthetic network, not a real one.}
Our network was generated by a computer program, not taken from
a real company. Real networks are messier --- they have legacy
systems, incomplete records, and irregular layouts. Our next
step is to test on an emulated network using Mininet and OpenFlow
(a real software-defined networking framework).

\textbf{Attack probabilities are static between updates.}
We treat the probability of each attack path as fixed. In reality
these change constantly. We handle updates manually via
Algorithm~\ref{alg:redeploy}, but we have not yet tested with
a live threat-intelligence feed.

\textbf{Detection is modelled as all-or-nothing.}
We assume a honeypot either catches an attacker 100\% of the
time or 0\% of the time. In reality, a sophisticated attacker
might recognise a honeypot and avoid it. A more realistic model
would use probabilities like 70\% or 80\% detection rates.

\textbf{Very large networks strain the solver.}
At XXLarge (5.5 million nodes), the final budget reconciliation
step adds noticeable delay. We plan to investigate faster
approximate solvers like CCLS and SATLike.

\subsection{What We Plan to Do Next}

\begin{enumerate}
  \item Build a real emulated testbed using Mininet with actual
        honeypot software (Cowrie for SSH, Conpot for industrial
        systems) to validate that our model's predictions match
        what happens in practice.
  \item Connect to STIX/TAXII threat-intelligence feeds so that
        our system automatically detects when attack probabilities
        change and triggers Algorithm~\ref{alg:redeploy}.
  \item Replace the binary detection model with a probability-based
        one, using real honeypot interaction logs to estimate
        detection rates per technique.
  \item Model a smarter attacker who can detect honeypots and
        change their route. This becomes a multi-round game where
        our solver updates the placement after each round.
\end{enumerate}

% ============================================================
\section{Conclusion}
\label{sec:conclusion}
% ============================================================

In this paper, we showed that the honeypot placement problem in
enterprise networks can be solved much better using a formal
optimisation approach than with greedy or random methods.

Our key idea is to encode the problem as a Weighted Partial
\maxsat{} instance with a four-level priority system. The highest
priority goes to catching attackers \emph{before} they reach their
target. Lower priorities handle technique coverage, tactic family
breadth, and basic detection opportunities. We also encode seven
types of hard constraints that must never be violated, including
budget limits, zone air gaps, and honeypot conflict rules.

The most important technical insight of our work is that
\emph{one honeypot can serve multiple attack paths at once}. A
greedy algorithm misses this because it evaluates one objective
at a time. Our solver encodes all attack paths simultaneously in
the same logical formula, and the \rcII{} algorithm automatically
discovers and exploits these overlaps.

Our experiments on six different network sizes, supported by
Monte Carlo attack simulations, showed consistent improvement over
all baselines. At the XLarge scale (500,000 computers), we achieved
$Q = 67.4$ versus the best baseline's $Q = 52.8$ --- a $+27.8\%$
improvement --- with twice the technique coverage and 60\% better
tactic-family breadth.

Most importantly, unlike any baseline method, our approach
\emph{proves} its solution is optimal. Defenders can deploy our
recommended placement with mathematical confidence that no other
deployment within their budget would perform better.

% ============================================================
\begin{thebibliography}{12}

\bibitem{hutchins2011intelligence}
E.~M.~Hutchins, M.~J.~Cloppert, and R.~M.~Amin,
``Intelligence-driven computer network defense informed by analysis
of adversary campaigns and intrusion kill chains,''
\textit{Leading Issues in Information Warfare \& Security Research},
vol.~1, no.~1, p.~80, 2011.

\bibitem{mitre2023}
MITRE Corporation,
``MITRE ATT\&CK Enterprise Matrix,'' 2023.
[Online]. Available: \url{https://attack.mitre.org/}

\bibitem{provos2004virtual}
N.~Provos,
``A virtual honeypot framework,''
in \textit{Proc.\ USENIX Security Symp.}, vol.~173, pp.~1--14, 2004.

\bibitem{spitzner2003honeypots}
L.~Spitzner,
\textit{Honeypots: Tracking Hackers}.
Addison-Wesley, 2003.

\bibitem{liffiton2016}
M.~Liffiton and A.~Malik,
``Enumerating infeasibility: Finding multiple MUSes quickly,''
in \textit{Proc.\ CPAIOR}, pp.~160--175, Springer, 2016.

\bibitem{narodytska2014}
N.~Narodytska and F.~Bacchus,
``Maximum satisfiability using core-guided MaxSAT resolution,''
in \textit{Proc.\ AAAI}, pp.~2717--2723, 2014.

\bibitem{tan2011}
G.~Tan, M.~Sherr, and W.~Zhou,
``Network-layer obliviousness: An attack-surface reduction technique,''
in \textit{Proc.\ ACM CCS}, 2011.

\bibitem{pawlick2019game}
J.~Pawlick, E.~Colbert, and Q.~Zhu,
``A game-theoretic taxonomy and survey of defensive deception for
cybersecurity and privacy,''
\textit{ACM Comput.\ Surv.}, vol.~52, no.~4, pp.~1--28, 2019.

\bibitem{applebaum2016intelligent}
A.~Applebaum \etal,
``Intelligent, automated red team emulation,''
in \textit{Proc.\ ACSAC}, 2016.

\bibitem{saha2015optimal}
S.~Saha and P.~Gera,
``Optimal placement of security sensors in enterprise networks:
An ILP approach,''
in \textit{Proc.\ IEEE ICC}, 2015.

\bibitem{jajodia2005topological}
S.~Jajodia, S.~Noel, and B.~O'Berry,
``Topological analysis of network attack vulnerability,''
in \textit{Managing Cyber Threats}, pp.~247--266, Springer, 2005.

\bibitem{zhu2013game}
Q.~Zhu and T.~Ba\c{s}ar,
``Game-theoretic approach to feedback-driven multi-stage moving
target defense,''
in \textit{Proc.\ GameSec}, pp.~246--263, Springer, 2013.

\end{thebibliography}

\end{document}
